\documentclass[11pt]{article}
\usepackage[margin=1.1in]{geometry}
\usepackage{amsmath,amssymb,amsthm}
\usepackage{booktabs}
\usepackage[colorlinks=true,linkcolor=blue,urlcolor=blue,citecolor=blue]{hyperref}

\newtheorem{theorem}{Theorem}[section]
\newtheorem{lemma}[theorem]{Lemma}
\newtheorem{corollary}[theorem]{Corollary}
\newtheorem{proposition}[theorem]{Proposition}
\newtheorem{conjecture}[theorem]{Conjecture}
\newtheorem{remark}[theorem]{Remark}

\newcommand{\qbin}[2]{\genfrac{[}{]}{0pt}{}{#1}{#2}}
\newcommand{\poch}[1]{(q;q)_{#1}}
\newcommand{\GK}{\mathrm{GK}}
\newcommand{\FERM}{\mathrm{FERM}}
\newcommand{\ZZ}{\mathbb{Z}}

\title{Two cases of the $\mathrm{A}_2$ Andrews--Gordon conjectures for
cylindric partitions ($k=3$ at modulus $11$, $k=4$ at modulus $13$),
and a bounded refinement of positivity at $d=4$}
\author{Robin Langer}
\date{July 2026}

\begin{document}
\maketitle

\begin{abstract}
Building on results of Corteel--Welsh, Kanade--Russell and Uncu, we prove
the first two cases beyond $k=2$ of Warnaar's fermionic-form conjectures
for cylindric partitions \cite[Conjectures~2.5 and~2.11]{W23}: the
balanced profiles $c=(3,3,2)$ at $d=8$ (modulus $11$, $k=3$) and
$c=(4,3,3)$ at $d=10$ (modulus $13$, $k=4$). Both proofs chain Warnaar's
own finite form \cite[Proposition~7.1]{W23}, in a coefficientwise limit,
with an exact Pochhammer split and one of Uncu's theorems. At modulus
$11$ one Kanade--Russell contiguous relation \cite[Lemma~9.2]{KR22} is
also needed to reach Uncu's modulus-$11$ theorem
\cite[Theorem~4.3]{Uncu23}; at modulus $13$ the chain is \emph{shorter}
--- the split lands verbatim on Uncu's modulus-$13$ theorem
\cite[Theorem~5.3]{Uncu23}, and no contiguous relation is required. As a
corollary, the positivity conjecture \cite[Conjecture~2.7]{W23} holds for
the orbits of $(3,3,2)$ and $(4,3,3)$. We also record a bounded refinement at $d=4$: for every profile
$c$ with $c_0+c_1+c_2=4$ one has
$H_{c,m}=\sum_{n\le m}\qbin{m}{n}Q_{n,c}$ with the $Q_{n,c}$ given by
explicit manifestly nonnegative polynomials built on the Corteel--Welsh
bounded forms; monotonicity of $(H_{c,m})_m$ follows. The positivity of
$Q_{n,c}$ at $d=4$ is Warnaar's result; the bounded expansion, and two
absorption identities that make the two ``wall'' orbits manifestly
positive, appear to be new. The elementary $q$-binomial infrastructure and
the $d=4$ argument have been formally verified in Lean~4/Mathlib.
\end{abstract}

\section{Introduction}\label{sec:intro}

Let $\GK_c(z,q)$ denote the generating function of cylindric partitions of
rank $3$ and profile $c=(c_0,c_1,c_2)$, in the notation of \cite{W23}, with
$z$ tracking the largest part; set $d:=c_0+c_1+c_2$ (the level), so the
associated modulus is $d+3$. Since $\GK_c$ is invariant under cyclic
rotation of the profile, every statement below about a profile applies
verbatim to its rotation orbit; we fix once and for all the profile
representatives used by Corteel--Welsh \cite{CW19} at $d=4$, namely
$(2,1,1)$, $(4,0,0)$, $(3,1,0)$, $(3,0,1)$, $(2,2,0)$, which lie in the
five distinct orbits.

Following \cite{W23}, for $\ell:=\gcd(d,3)$ define
\begin{equation}\label{eq:Qdef}
Q_{n,c}(q):=(q^{\ell};q^{\ell})_n\,[z^n]\Bigl((zq;q)_\infty\,
\GK_c(z,q)\Bigr).
\end{equation}
The positivity conjecture, stated by Corteel--Dousse--Uncu
\cite[Conjecture~4.2]{CDU22} and reproduced as
\cite[Conjecture~2.7]{W23}, reads:

\begin{conjecture}[{\cite[Conjecture~2.7]{W23}}]\label{con:pos}
Let $c=(c_0,c_1,c_2)$ and $d:=c_0+c_1+c_2$ such that
$d\not\equiv 0\pmod 3$. Then $Q_{n,(c_0,c_1,c_2)}(q)$ is a polynomial in
$q$ with nonnegative coefficients. Moreover,
$Q_{n,(c_0,c_1,c_2)}(1)=\bigl(\tfrac16(d+1)(d+2)-1\bigr)^n$.
\end{conjecture}

Polynomiality and the evaluation at $q=1$ were proved by Welsh
\cite{Welsh21}; Warnaar proved positivity for $k=1$ and $k=2$, i.e.\ for
all profiles with $d\in\{2,4,5\}$ \cite{W23}; positivity for general $d$
remains open. For the balanced profile $(k,k,k-1)$, $d=3k-1$, Warnaar
conjectured an explicit fermionic form:

\begin{conjecture}[{\cite[Conjecture~2.5]{W23}, first identity}]
\label{con:ferm}
For a positive integer $k$,
\[
\GK_{(k,k,k-1)}(z,q)
=\frac{1}{(zq;q)_\infty}
\sum_{\substack{n_1,\dots,n_k\ge 0\\ m_1,\dots,m_{k-1}\ge 0}}
\frac{z^{n_1}q^{n_k^2}}{\poch{n_1}}
\prod_{i=1}^{k-1}q^{n_i^2-n_im_i+m_i^2}
\qbin{n_i}{n_{i+1}}\qbin{n_i-n_{i+1}+m_{i+1}}{m_i},
\]
where $m_k:=2n_k$.
\end{conjecture}

The modulus-$(3k+1)$ companion concerns the balanced profile
$(k,k-1,k-1)$ at level $d=3k-2$:

\begin{conjecture}[{\cite[Conjecture~2.11]{W23}, first identity}]
\label{con:fermb}
For a positive integer $k$,
\begin{multline*}
\GK_{(k,k-1,k-1)}(z,q)
=\frac{1}{(zq;q)_\infty}
\sum_{\substack{n_1,\dots,n_{k-1}\ge 0\\ m_1,\dots,m_{k-1}\ge 0}}
\frac{z^{n_1}q^{\sum_{i=1}^{k-1}(n_i^2-n_im_i+m_i^2)}}{\poch{n_1}}\\
\times\qbin{2n_{k-1}}{m_{k-1}}
\prod_{i=1}^{k-2}\qbin{n_i}{n_{i+1}}\qbin{n_i-n_{i+1}+m_{i+1}}{m_i}.
\end{multline*}
\end{conjecture}

Warnaar proved Conjecture~\ref{con:ferm} for $k=1$ and $k=2$
\cite[Theorem~2.6]{W23}; the $k=1$ case of Conjecture~\ref{con:fermb} is
trivial and $k=2$ is due to Corteel--Welsh \cite[Theorem~3.2]{CW19}. As
Warnaar observes, either conjecture at general $k$ immediately implies
manifestly positive representations for the $Q_{n,c}$ of its balanced
profile, hence Conjecture~\ref{con:pos} there.

The purpose of this note is threefold.

\emph{First} (Section~\ref{sec:equiv}), we record a simple but, we
believe, clarifying observation: for $\gcd(d,3)=1$, writing
$F_{c,m}(q)$ for the generating function of
cylindric partitions of profile $c$ with largest part at most $m$, and
$H_{c,m}:=\poch{m}F_{c,m}$, the family $(H_{c,m})_m$ has a
\emph{unique} expansion $H_{c,m}=\sum_{n\le m}\qbin{m}{n}a_n$ with
$a_n$ independent of $m$, and necessarily $a_n=Q_{n,c}$. Consequently
Conjecture~\ref{con:pos} at $c$ is \emph{equivalent} to the existence of a
coefficientwise-nonnegative $q$-binomial expansion of $(H_{c,m})_m$. This
is a precise sense in which the problem of bounded $\mathrm{A}_2$
Andrews--Gordon identities and the positivity conjecture are two faces of
one problem.

\emph{Second} (Section~\ref{sec:d8}), we prove
Conjecture~\ref{con:ferm} for $k=3$ ($d=8$, modulus $11$) and
Conjecture~\ref{con:fermb} for $k=4$ ($d=10$, modulus $13$). Each proof
is a short chain of results from the literature: Warnaar's finite form
\cite[Proposition~7.1]{W23} in a coefficientwise limit, an exact
single-step Pochhammer split, and one of Uncu's theorems
\cite[Theorems~4.3 and~5.3]{Uncu23}; the modulus-$11$ chain additionally
needs one contiguous relation of Kanade--Russell \cite[Lemma~9.2]{KR22},
while the modulus-$13$ chain --- although it concerns the larger $k$ ---
is strictly shorter and needs none. We make no claim of depth here ---
all the hard analysis is in \cite{W23,KR22,Uncu23} --- but the resulting
statements appear to be new, and they settle the first cases of
Conjecture~\ref{con:pos} at $d=8$ and $d=10$.

\emph{Third} (Section~\ref{sec:d4}), we revisit $d=4$, where positivity is
Warnaar's result. Starting from the bounded generating functions of
Corteel--Welsh \cite[Theorem~1.2]{CW19} we obtain, for all five orbits
including the two ``wall'' orbits $(3,0,1)$ and $(2,2,0)$, the bounded
fermionic-form expansion $H_{c,m}=\sum_{n\le m}\qbin{m}{n}Q_{n,c}$
together with explicit manifestly nonnegative formulas for the $Q_{n,c}$.
The wall orbits require two absorption identities
(Lemmas~\ref{lem:A} and~\ref{lem:B}), proved by elementary $q$-Pascal
manipulations; Lemma~\ref{lem:B} is an exact identity, not merely an
inequality. Monotonicity $H_{c,m}\ge H_{c,m-1}$ follows for all orbits.

Section~\ref{sec:lean} describes a formal verification: the $q$-binomial
infrastructure of Section~\ref{sec:equiv}, both absorption lemmas, the
positivity of all five $d=4$ forms, and the assembly logic of
Section~\ref{sec:d8} have been proved in Lean~4 over Mathlib, with the
literature inputs entering as explicitly named hypotheses. We also report
exact $\ZZ[q]$ verification of Conjecture~\ref{con:pos} across fourteen
levels up to $d=31$.

Finally, Appendix~\ref{app:method} describes how this note was produced:
the mathematics was generated and checked by an orchestrated system of AI
agents, with the author acting as guarantor of the results.

\section{The $q$-binomial expansion of the bounded family}
\label{sec:equiv}

Throughout this section fix a profile $c$ with $\ell=\gcd(d,3)=1$, and
work in $\ZZ[[q]]$. Let $g_m$ be the generating function of cylindric
partitions of profile $c$ whose largest part is exactly $m$ (so $g_0=1$),
let $F_{c,m}=\sum_{j\le m}g_j$, and set
\[
H_m:=\poch{m}F_{c,m},\qquad
Q_n:=Q_{n,c}(q)=\poch{n}\,[z^n]\Bigl((zq;q)_\infty
\textstyle\sum_{m\ge0}g_mz^m\Bigr).
\]

\begin{lemma}[truncated Euler telescoping]\label{lem:E}
For every integer $K\ge 0$,
\[
\sum_{k=0}^{K}\frac{(-1)^k q^{\binom{k}{2}}}{\poch{k}}
=\frac{(-1)^K q^{\binom{K+1}{2}}}{\poch{K}} .
\]
\end{lemma}

\begin{proof}
Induction on $K$: both sides are $1$ at $K=0$, and adding the
$(K{+}1)$-st term to the right side gives
$(-1)^{K+1}q^{\binom{K+1}{2}}\bigl(1-(1-q^{K+1})\bigr)/\poch{K+1}
=(-1)^{K+1}q^{\binom{K+2}{2}}/\poch{K+1}$.
\end{proof}

\begin{theorem}[the $Q$-transform]\label{thm:Q}
For every $n\ge 0$,
$\displaystyle
Q_n=\sum_{m=0}^{n}(-1)^{n-m}q^{\binom{n-m}{2}}\qbin{n}{m}H_m$.
\end{theorem}

\begin{proof}
By Euler's expansion $(x;q)_\infty=\sum_{k\ge0}(-1)^kq^{\binom{k}{2}}
x^k/\poch{k}$ at $x=zq$, extracting the coefficient of $z^n$,
\begin{equation}\label{eq:Qeuler}
Q_n=\poch{n}\sum_{m=0}^{n}
\frac{(-1)^{n-m}q^{\binom{n-m+1}{2}}}{\poch{n-m}}\,g_m .
\end{equation}
Substituting $H_m=\poch{m}\sum_{j\le m}g_j$ into the claimed sum, using
$\qbin{n}{m}\poch{m}=\poch{n}/\poch{n-m}$, interchanging the (finite)
sums, and applying Lemma~\ref{lem:E} to the inner sum over $k=n-m$
reproduces \eqref{eq:Qeuler}.
\end{proof}

\begin{theorem}[Gauss inversion; classical]\label{thm:G}
For sequences $(a_n),(b_n)$ in $\ZZ[[q]]$ the statements
\textup{(i)} $b_n=\sum_{j\le n}\qbin{n}{j}a_j$ for all $n$ and
\textup{(ii)} $a_n=\sum_{m\le n}(-1)^{n-m}q^{\binom{n-m}{2}}\qbin{n}{m}b_m$
for all $n$ are equivalent; moreover for given $(b_n)$ the sequence
$(a_n)$ in \textup{(i)} is unique.
\end{theorem}

\begin{proof}
Both transforms are lower-unitriangular, so it suffices to check the
orthogonality
\[
\sum_{m=k}^{n}(-1)^{n-m}q^{\binom{n-m}{2}}\qbin{n}{m}\qbin{m}{k}
=\delta_{n,k} .
\]
Using $\qbin{n}{m}\qbin{m}{k}=\qbin{n}{k}\qbin{n-k}{m-k}$ and
substituting $i=n-m$, $J=n-k$, the left side equals
\[
\qbin{n}{k}\sum_{i=0}^{J}(-1)^iq^{\binom{i}{2}}\qbin{J}{i}
=\qbin{n}{k}\,(x;q)_J\Big|_{x=1}=\qbin{n}{k}\,\delta_{J,0}
\]
by Gauss's finite $q$-binomial theorem.
\end{proof}

\begin{corollary}[inverse $Q$-transform; unconditional]\label{cor:inv}
For every $m\ge0$,
$\displaystyle H_m=\sum_{n=0}^{m}\qbin{m}{n}Q_n$.
\end{corollary}

At $q=1$ this reconciles Welsh's evaluation with the orbit count:
$H_m(1)=\sum_n\binom{m}{n}(K-1)^n=K^m$, $K=\tfrac16(d+1)(d+2)$.

\begin{proposition}[equivalence]\label{prop:equiv}
Fix $c$ with $d\not\equiv0\pmod3$.
\begin{enumerate}
\item[\textup{(1)}] If $H_m=\sum_{j\le m}\qbin{m}{j}a_j$ for all $m$ with
$a_j\in\ZZ[[q]]$ independent of $m$, then $a_j=Q_j$ for all $j$. In
particular, any bounded fermionic form for $(H_m)_m$ --- an expansion of
this shape with manifestly nonnegative $a_j$ --- proves
Conjecture~\ref{con:pos} at $c$, with the explicit formula $Q_n=a_n$.
\item[\textup{(2)}] The family $(H_m)_m$ admits a
coefficientwise-nonnegative $q$-binomial expansion \emph{if and only if}
Conjecture~\ref{con:pos} holds at $c$.
\end{enumerate}
\end{proposition}

\begin{proof}
(1) is the uniqueness in Theorem~\ref{thm:G} combined with
Corollary~\ref{cor:inv}. (2): the unique expansion coefficients are the
$Q_n$, so a nonnegative expansion exists iff they are all nonnegative.
\end{proof}

Thus, for $\gcd(d,3)=1$, there is no shortcut criterion certifying the
existence of a nonnegative expansion without exhibiting its coefficients:
for the family $(H_m)_m$, existence \emph{is} the conjecture.

\section{Conjectures~\ref{con:ferm} and~\ref{con:fermb} beyond $k=2$:
the profiles $(3,3,2)$ and $(4,3,3)$}\label{sec:d8}

Write $\FERM_k(z,q)$ for the multisum in Conjecture~\ref{con:ferm}
(without the $1/(zq;q)_\infty$ prefactor), and
$G_c(z,q):=(zq;q)_\infty\GK_c(z,q)$,
$\mathcal{H}_c(z,q):=G_c(z,q)/\poch{\infty}$. For $\rho,\sigma\in\ZZ^3$
let, following Uncu \cite[Eq.~(2.8)]{Uncu23} (the case of modulus $11$),
abbreviating
\[
P_{r,s}:=\poch{r_1-r_2}\poch{r_2-r_3}\poch{r_3}\,
\poch{s_1-s_2}\poch{s_2-s_3}\poch{s_3},
\]
\begin{equation}\label{eq:S11}
S_{11}(z;\rho\,|\,\sigma):=
\sum_{\substack{r_1\ge r_2\ge r_3\ge 0\\ s_1\ge s_2\ge s_3\ge 0}}
z^{r_1}\,
\frac{q^{\sum_{i=1}^{3}(r_i^2-r_is_i+s_i^2+\rho_ir_i+\sigma_is_i)+2r_3s_3}}
{P_{r,s}\,\poch{r_3+s_3+1}} ,
\end{equation}
and abbreviate $e_3=(0,0,0)$, $e_2=(0,0,1)$, $e_1=(0,1,1)$, with
$\delta_i$ the $i$-th unit vector. We suppress $z$ and write
$S(\rho\,|\,\sigma)$.

\begin{theorem}\label{thm:d8}
$G_{(3,3,2)}(z,q)=\FERM_3(z,q)$; that is, Conjecture~\ref{con:ferm} holds
for $k=3$.
\end{theorem}

\begin{corollary}\label{cor:d8pos}
For all $n\ge 0$,
\[
Q_{n,(3,3,2)}(q)=
\sum_{\substack{n\ge n_2\ge n_3\ge 0\\ m_1,m_2\ge 0}}
q^{\,n^2+n_2^2+n_3^2-nm_1-n_2m_2+m_1^2+m_2^2}
\qbin{n}{n_2}\qbin{n-n_2+m_2}{m_1}\qbin{n_2}{n_3}\qbin{n_2+n_3}{m_2},
\]
a manifestly nonnegative polynomial. In particular
Conjecture~\ref{con:pos} holds for the orbit of $(3,3,2)$ at $d=8$.
\end{corollary}

\begin{proof}[Proof of the corollary]
$Q_{n,(3,3,2)}=\poch{n}[z^n]\FERM_3$; the summand at $n_1=n$ carries the
denominator $\poch{n}$, which cancels. (Here
$n_2-n_3+m_3=n_2+n_3$ since $m_3=2n_3$, and the exponent $n_3^2$ has been
absorbed into the displayed quadratic form.)
\end{proof}

\begin{proof}[Proof of Theorem~\ref{thm:d8}]
The proof chains four proved statements.

\emph{Step 1: Warnaar's finite form in the limit.}
For $\sigma=(1,\dots,1,a)$, $a\in\{-1,1\}$, Warnaar defines
\cite[Eqs.~(7.1),(7.2)]{W23}
\[
F^{(a)}_{n_0,m_0;k}(z,q)=
\sum_{\substack{n_1,\dots,n_k\ge 0\\ m_1,\dots,m_k\ge 0}}
\frac{z^{n_1}}{\poch{n_k+m_k}}
\prod_{i=1}^{k}q^{n_i^2-\sigma_in_im_i+m_i^2}
\qbin{n_{i-1}}{n_i}\qbin{m_{i-1}}{m_i},
\]
and proves \cite[Proposition~7.1, case $a=-1$]{W23} the polynomial
identity
\begin{equation}\label{eq:ff}
F^{(-1)}_{n_0,m_0;k}(z,q)=
\sum_{\substack{n_1,\dots,n_k\ge 0\\ m_1,\dots,m_{k-1}\ge 0}}
\frac{z^{n_1}q^{n_k^2}}{\poch{m_0-n_1+m_1}}\qbin{n_0}{n_1}
\prod_{i=1}^{k-1}q^{n_i^2-n_im_i+m_i^2}
\qbin{n_i}{n_{i+1}}\qbin{n_i-n_{i+1}+m_{i+1}}{m_i},
\end{equation}
where $m_k:=2n_k$. Fix $k=3$ and let $n_0,m_0\to\infty$. Each coefficient
of $z^nq^N$ on both sides stabilizes; the right side of \eqref{eq:ff}
becomes $\FERM_3(z,q)/\poch{\infty}$, while on the left the two binomial
chains telescope,
$\qbin{n_0}{n_1}\qbin{n_1}{n_2}\qbin{n_2}{n_3}\to
1/(\poch{n_1-n_2}\poch{n_2-n_3}\poch{n_3})$ and likewise for the
$m$-chain. Renaming $(n,m)\to(r,s)$ and using $\sigma=(1,1,-1)$, so that
the $i=3$ exponent is $r_3^2+r_3s_3+s_3^2=(r_3^2-r_3s_3+s_3^2)+2r_3s_3$,
we obtain
\begin{equation}\label{eq:T}
\frac{\FERM_3(z,q)}{\poch{\infty}}=T(z,q):=
\sum_{\substack{r_1\ge r_2\ge r_3\ge 0\\ s_1\ge s_2\ge s_3\ge 0}}
z^{r_1}\,
\frac{q^{\sum_{i=1}^{3}(r_i^2-r_is_i+s_i^2)+2r_3s_3}}
{P_{r,s}\,\poch{r_3+s_3}} .
\end{equation}

\emph{Step 2: Pochhammer split.}
Using $1/\poch{r_3+s_3}=(1-q^{r_3+s_3+1})/\poch{r_3+s_3+1}$ termwise,
together with $q^{r_3+s_3+1}=q\cdot q^{r_3}\cdot q^{s_3}$ --- which
realizes the subtracted term as the shift
$(\rho_3,\sigma_3)\mapsto(\rho_3+1,\sigma_3+1)$ in \eqref{eq:S11} with a
global factor $q$ --- we get
\begin{equation}\label{eq:split}
T(z,q)=S(e_3\,|\,e_3)-q\,S(e_2\,|\,e_2).
\end{equation}

\emph{Step 3: one Kanade--Russell contiguous relation.}
Kanade--Russell \cite[Lemma~9.2]{KR22} prove, for modulus
$\equiv-1\pmod 3$ and $\sigma_{k-1}=0$ (in the notation
\eqref{eq:S11}, $\sigma_2=0$), the relation
\[
S(\rho\,|\,\sigma)-S(\rho\,|\,\sigma+\delta_3)
-q\,S(\rho+\delta_3\,|\,\sigma+\delta_3)
+q\,S(\rho+\delta_3\,|\,\sigma+\delta_2+\delta_3)=0 .
\]
At $\rho=\sigma=(0,0,0)$ this reads
$S(e_3|e_3)-S(e_3|e_2)-qS(e_2|e_2)+qS(e_2|e_1)=0$, i.e.
\begin{equation}\label{eq:bridge}
S(e_3\,|\,e_3)-q\,S(e_2\,|\,e_2)=S(e_3\,|\,e_2)-q\,S(e_2\,|\,e_1).
\end{equation}

\emph{Step 4: Uncu's modulus-$11$ theorem.}
Uncu \cite[Eq.~(4.1) and Theorem~4.3]{Uncu23} proves
\begin{equation}\label{eq:uncu}
\mathcal{H}_{(3,3,2)}(z,q)=S(e_3\,|\,e_2)-q\,S(e_2\,|\,e_1).
\end{equation}

Combining \eqref{eq:T}--\eqref{eq:uncu},
$\FERM_3(z,q)=\poch{\infty}\,\mathcal{H}_{(3,3,2)}(z,q)
=G_{(3,3,2)}(z,q)$.
\end{proof}

\subsection*{The profile $(4,3,3)$ at modulus $13$}

Write $\FERM^{+}_{k-1}(z,q)$ for the multisum in
Conjecture~\ref{con:fermb} (without the $1/(zq;q)_\infty$ prefactor).
Following Uncu \cite[Eq.~(2.10)]{Uncu23} (the case of modulus $13$),
define
\begin{equation}\label{eq:S13}
S_{13}(z;\rho\,|\,\sigma):=
\sum_{\substack{r_1\ge r_2\ge r_3\ge 0\\ s_1\ge s_2\ge s_3\ge 0}}
z^{r_1}\,
\frac{q^{\sum_{i=1}^{3}(r_i^2-r_is_i+s_i^2+\rho_ir_i+\sigma_is_i)}}
{P_{r,s}\,\poch{r_3+s_3+1}} ;
\end{equation}
note that, in contrast to \eqref{eq:S11}, there is \emph{no} cross term
$q^{2r_3s_3}$ in the modulus-$13$ family. We again suppress $z$.

\begin{theorem}\label{thm:d10}
$G_{(4,3,3)}(z,q)=\FERM^{+}_3(z,q)$; that is, Conjecture~\ref{con:fermb}
holds for $k=4$.
\end{theorem}

\begin{corollary}\label{cor:d10pos}
For all $n\ge 0$,
\begin{multline*}
Q_{n,(4,3,3)}(q)=
\sum_{\substack{n_2,n_3\ge 0\\ m_1,m_2,m_3\ge 0}}
q^{\,n^2+n_2^2+n_3^2-nm_1-n_2m_2-n_3m_3+m_1^2+m_2^2+m_3^2}\\
\times\qbin{n}{n_2}\qbin{n-n_2+m_2}{m_1}
\qbin{n_2}{n_3}\qbin{n_2-n_3+m_3}{m_2}\qbin{2n_3}{m_3},
\end{multline*}
a manifestly nonnegative polynomial. In particular
Conjecture~\ref{con:pos} holds for the orbit of $(4,3,3)$ at $d=10$.
\end{corollary}

\begin{proof}[Proof of Theorem~\ref{thm:d10} (sketch)]
The chain has three links --- Steps 1, 2 and 4 above have exact
analogues, and there is no analogue of Step 3. A detailed writeup is in
the companion note \cite{Comp26}.

\emph{Step $1'$.} The $a=+1$ case of \cite[Proposition~7.1]{W23} at
$k=3$ (all $\sigma_i=1$), in the same coefficientwise limit
$n_0,m_0\to\infty$ as in Step~1, gives
\[
\frac{\FERM^{+}_3(z,q)}{\poch{\infty}}
=\sum_{\substack{r_1\ge r_2\ge r_3\ge 0\\ s_1\ge s_2\ge s_3\ge 0}}
z^{r_1}\,
\frac{q^{\sum_{i=1}^{3}(r_i^2-r_is_i+s_i^2)}}{P_{r,s}\,\poch{r_3+s_3}}
=:T'(z,q),
\]
with no cross term, since $\sigma_3=+1$ leaves the quadratic form
$r_3^2-r_3s_3+s_3^2$ intact.

\emph{Step $2'$.} The identical Pochhammer split
$1/\poch{r_3+s_3}=(1-q^{r_3+s_3+1})/\poch{r_3+s_3+1}$ gives, termwise,
$T'=S_{13}(e_3\,|\,e_3)-q\,S_{13}(e_2\,|\,e_2)$.

\emph{Step $4'$.} Uncu \cite[Eq.~(5.1), last entry, and
Theorem~5.3]{Uncu23} proves
$\mathcal{H}_{(4,3,3)}(z,q)=S_{13}(e_3\,|\,e_3)-q\,S_{13}(e_2\,|\,e_2)$
--- \emph{verbatim} the right side of Step $2'$. Hence
$\FERM^{+}_3=\poch{\infty}\mathcal{H}_{(4,3,3)}=G_{(4,3,3)}$.
\end{proof}

\begin{remark}[why the $d=10$ chain is shorter]\label{rem:scale}
At modulus $3k-1$ the finite form produces the pair
$(e_3|e_3),(e_2|e_2)$ while the proved Kanade--Russell expression lists
$(e_3|e_2),(e_2|e_1)$, so one contiguous relation is needed to bridge
them (Step~3). At modulus $3k+1$ the Kanade--Russell pattern
$\mathcal{H}_{(c_1,c_2,c_3)}=S(e_{c_2}|e_{c_3})-qS(e_{c_2-1}|e_{c_3-1})$,
evaluated at the balanced profile $c_2=c_3=k-1$, is exactly the split
pair, and no bridge is needed. Thus the template --- finite form,
coefficientwise limit, Pochhammer split, Uncu --- scales directly to
whichever moduli have proved Kanade--Russell expressions, with the
contiguous-relation step required only at moduli $\equiv-1\pmod 3$.
\end{remark}

\begin{remark}[provenance of Uncu's theorems]\label{rem:prov}
Uncu's Theorems~4.3 and~5.3 \cite{Uncu23} are computer-assisted proofs:
they rest on ideal-membership computations (Gaussian elimination over
$q$-difference operators) whose certificates are provided in the
ancillary files of arXiv:2301.01359. Theorems~\ref{thm:d8}
and~\ref{thm:d10} inherit that dependency.
\end{remark}

\section{$d=4$: bounded fermionic forms and the absorption identities}
\label{sec:d4}

Fix $d=4$ (modulus $7$), and recall the five orbit representatives
$(2,1,1)$, $(4,0,0)$, $(3,1,0)$, $(3,0,1)$, $(2,2,0)$ of
Section~\ref{sec:intro}. Write $T(n,j):=q^{n^2+j^2-nj}$. Corteel--Welsh
\cite[Theorem~1.2]{CW19} give the bounded generating functions
\[
F_{c,m}=\sum_{n\le m}\frac{\widetilde A_{n,c}}{\poch{m-n}},
\]
with $m$-independent numerators $\widetilde A_{n,c}$; explicitly, reading
off from their theorem (equivalently from the $G_c$ forms of
\cite[Theorem~3.2]{CW19}),
\begin{align*}
\widetilde A_n^{(2,1,1)}&=\sum_{j=0}^{2n}\frac{T(n,j)}{\poch{n}}
\qbin{2n}{j},\qquad
\widetilde A_n^{(4,0,0)}=\sum_{j}\frac{T(n,j)\,q^{n+j}}{\poch{n}}
\qbin{2n}{j},\\
\widetilde A_n^{(3,1,0)}&=\sum_{j}\frac{T(n,j)\,q^{j}}{\poch{n}}
\qbin{2n}{j},\\
\widetilde A_n^{(3,0,1)}&=\sum_{j}T(n,j)\Bigl(
\frac{q^{n}}{\poch{n}}\qbin{2n}{j}
+\frac{q^{2j}}{\poch{n-1}}\qbin{2n-2}{j}\Bigr),\\
\widetilde A_n^{(2,2,0)}&=\sum_{j}T(n,j)\Bigl(
\frac{q^{n}}{\poch{n}}\qbin{2n}{j}
+\frac{q^{j}(1+q^{n+j})}{\poch{n-1}}\qbin{2n-2}{j}\Bigr),
\end{align*}
with the convention $1/\poch{-1}=0$.

Since the $\widetilde A_{n,c}$ are independent of $m$, letting
$m\to\infty$ gives $Q_{n,c}=\poch{n}\widetilde A_{n,c}$, and
Proposition~\ref{prop:equiv}(1) (or a one-line direct computation with
$\qbin{m}{n}=\poch{m}/(\poch{m-n}\poch{n})$) yields
\begin{equation}\label{eq:bffd4}
H_{c,m}=\sum_{n=0}^{m}\qbin{m}{n}\,Q_{n,c}
\qquad\text{for all five orbits.}
\end{equation}
Using $\poch{n}/\poch{n-1}=1-q^n$, the coefficients are
\begin{align}
Q_n^{(2,1,1)}&=\sum_{j}T(n,j)\qbin{2n}{j},\qquad
Q_n^{(4,0,0)}=\sum_{j}T(n,j)\,q^{n+j}\qbin{2n}{j},\qquad
Q_n^{(3,1,0)}=\sum_{j}T(n,j)\,q^{j}\qbin{2n}{j},
\notag\\
Q_n^{(3,0,1)}&=X_n+(1-q^n)X'_n,\qquad
Q_n^{(2,2,0)}=X_n+(1-q^n)Y'_n,
\label{eq:wallforms}
\end{align}
where
\begin{gather*}
X_n=\sum_jT(n,j)\,q^n\qbin{2n}{j},\qquad
X'_n=\sum_jT(n,j)\,q^{2j}\qbin{2n-2}{j},\\
Y'_n=\sum_jT(n,j)\,q^{j}(1+q^{n+j})\qbin{2n-2}{j}.
\end{gather*}
The first three forms are manifestly nonnegative. The two wall orbits
\eqref{eq:wallforms} are not, because of the factor $1-q^n$; positivity
follows from the absorption inequalities $X_n\ge q^nX'_n$ and
$X_n\ge q^nY'_n$, which we now prove.

\begin{lemma}[Absorption A]\label{lem:A}
For all $n\ge1$,
\[
X_n-q^nX'_n=\sum_jT(n,j)\,q^n\Bigl((q^{j-1}+q^{j})\qbin{2n-2}{j-1}
+\qbin{2n-2}{j-2}\Bigr)\;\ge\;0 .
\]
\end{lemma}

\begin{proof}
Two applications of the $q$-Pascal rule
$\qbin{N}{j}=\qbin{N-1}{j-1}+q^j\qbin{N-1}{j}$ give
\[
\qbin{2n}{j}=q^{2j}\qbin{2n-2}{j}
+(q^{j-1}+q^{j})\qbin{2n-2}{j-1}+\qbin{2n-2}{j-2}.
\]
The first term contributes exactly $q^n$ times the $j$-th summand of
$X'_n$; the remainder is nonnegative.
\end{proof}

\begin{lemma}[Absorption B; exact]\label{lem:B}
For all $n\ge1$,
\[
X_n-q^nY'_n
=\sum_{j=0}^{2n-1}q^{\,n^2+j^2-nj+2j+1}\qbin{2n-1}{j}\;\ge\;0 .
\]
\end{lemma}

\begin{proof}
Divide by the global factor $q^{n^2+n}$ and set
\[
D_n=\sum_jq^{j^2-nj}\qbin{2n}{j}
-\sum_jq^{j^2-nj}\bigl(q^{j}+q^{n+2j}\bigr)\qbin{2n-2}{j}.
\]
\emph{Step 1.} The two $q$-Pascal rules
$\qbin{2n}{j}=\qbin{2n-1}{j-1}+q^j\qbin{2n-1}{j}$ and
$\qbin{2n-1}{j}=\qbin{2n-2}{j}+q^{2n-1-j}\qbin{2n-2}{j-1}$ combine to
\begin{equation}\label{eq:star}
\qbin{2n}{j}-q^{j}\qbin{2n-2}{j}
=\qbin{2n-1}{j-1}+q^{2n-1}\qbin{2n-2}{j-1},
\end{equation}
valid for all $0\le j\le 2n$ with $\qbin{N}{k}:=0$ for $k<0$ or $k>N$.

\emph{Step 2.} Substituting \eqref{eq:star},
\[
D_n=\sum_jq^{j^2-nj}\qbin{2n-1}{j-1}
+\sum_jq^{j^2-nj+2n-1}\qbin{2n-2}{j-1}
-\sum_{j=0}^{2n-2}q^{j^2-nj+n+2j}\qbin{2n-2}{j}.
\]
In the last sum replace $j\mapsto j-1$: the exponent becomes
$(j-1)^2-n(j-1)+n+2(j-1)=j^2-nj+2n-1$, so it cancels the middle sum
\emph{exactly}. Hence $D_n=\sum_{j\ge1}q^{j^2-nj}\qbin{2n-1}{j-1}$;
re-indexing $j\mapsto j+1$ and restoring $q^{n^2+n}$ gives the claim.
\end{proof}

\begin{theorem}\label{thm:d4}
For every profile $c$ with $c_0+c_1+c_2=4$:
\begin{enumerate}
\item[\textup{(i)}] $Q_{n,c}(q)\ge0$ coefficientwise for all $n$, with
the manifestly positive formulas above (walls via
Lemmas~\ref{lem:A},~\ref{lem:B}); $Q_{n,c}(1)=4^n$.
\item[\textup{(ii)}] \textup{(bounded fermionic form at $d=4$)}
$H_{c,m}=\sum_{n\le m}\qbin{m}{n}Q_{n,c}$ with $Q_{n,c}\ge0$, for all
five orbits including the walls.
\item[\textup{(iii)}] \textup{(monotonicity)}
$H_{c,m}-H_{c,m-1}
=\sum_{n\ge1}q^{m-n}\qbin{m-1}{n-1}Q_{n,c}\;\ge\;0$.
\end{enumerate}
\end{theorem}

\begin{proof}
(i): the forms \eqref{eq:wallforms} plus Lemmas~\ref{lem:A}
and~\ref{lem:B}. (ii): equation \eqref{eq:bffd4}. (iii): the $q$-Pascal
rule $\qbin{m}{n}-\qbin{m-1}{n}=q^{m-n}\qbin{m-1}{n-1}$ applied termwise
to (ii).
\end{proof}

We emphasize that part~(i) is not new: positivity of $Q_{n,c}$ for all
profiles with $d\in\{2,4,5\}$ is Warnaar's result \cite{W23}. What we
believe to be worth recording is the bounded refinement (ii)--(iii) ---
by Proposition~\ref{prop:equiv} the expansion \eqref{eq:bffd4} is the
\emph{unique} bounded fermionic form at the first level --- and the
absorption mechanism itself: the wall orbits satisfy
$Q_n=P_n+(1-q^n)P'_n$ with $P_n,P'_n\ge0$, and positivity reduces to an
inequality $P_n\ge q^nP'_n$ provable purely by $q$-Pascal manipulation
(a double-Pascal split for one wall, an exact shift-cancellation for the
other). It seems plausible that this pattern persists for the
non-balanced orbits at higher $d$, where positive systems of
Corteel--Welsh type are not yet available.

\section{Machine verification}\label{sec:lean}

The elementary layers of Sections~\ref{sec:equiv}--\ref{sec:d4} have been
formalized in Lean~4 over Mathlib; the development (library
\texttt{WarnaarGlue}, five modules, no \texttt{sorry}) is available at
\begin{center}
\url{https://github.com/RaggedR/warnaar-glue}
\end{center}
$q$-binomial coefficients are defined by the Pascal recursion over an
arbitrary commutative ring (specialized to
$\ZZ[q]=\texttt{Polynomial}\;\ZZ$ where positivity is involved), and
nonnegativity is the predicate ``every coefficient of the polynomial is
$\ge0$''. Table~\ref{tab:lean} maps the statements of this note to the
formal theorems.

\begin{table}[ht]
\centering
\small
\begin{tabular}{ll}
\toprule
Statement in this note & Lean theorem(s) \\
\midrule
$q$-Pascal rules & \texttt{pascal}$_1$, \texttt{pascal}$_2$ \\
Gauss's finite $q$-binomial theorem at $x=1$ & \texttt{alt\_sum} \\
orthogonality in Theorem~\ref{thm:G} & \texttt{orth\_ML},
\texttt{orth\_LM} \\
Theorem~\ref{thm:G} (Gauss inversion, both directions) &
\texttt{qbinom\_inversion} \\
Theorem~\ref{thm:Q} $\leftrightarrow$ Corollary~\ref{cor:inv}
(abstract form) & \texttt{Q\_transform\_of\_H},
\texttt{H\_of\_Q\_transform} \\
monotonicity kernel (Theorem~\ref{thm:d4}(iii)) &
\texttt{pascal\_ladder} \\
Pochhammer split kernel of Step 2, Section~\ref{sec:d8} &
\texttt{qpoch\_split}, \texttt{qpoch\_split\_shift} \\
assembly of Steps 1--4 ($d=8$ chain) &
\texttt{seed2\_assembly}, \texttt{seed2\_chain} \\
$d=10$ chain (Theorem~\ref{thm:d10}) & not formalized
(cf.\ \texttt{seed2\_chain}) \\
Lemma~\ref{lem:A} & \texttt{absorption\_A} \\
Lemma~\ref{lem:B} & \texttt{absorption\_B} \\
positivity of all five $d=4$ forms & \texttt{Qcw\_nonneg} \\
identification $Q_{n,c}=\poch{n}\widetilde A_{n,c}$ via inversion &
\texttt{d4\_Q\_eq\_Qcw} \\
Theorem~\ref{thm:d4}(i) & \texttt{d4\_positive} \\
Theorem~\ref{thm:d4}(ii) & \texttt{d4\_BFF} \\
Theorem~\ref{thm:d4}(iii) & \texttt{d4\_monotone} \\
\bottomrule
\end{tabular}
\caption{Paper-to-Lean dictionary.}
\label{tab:lean}
\end{table}

The division of labour is made explicit in the formalization: results
proved here from scratch (the $q$-binomial identities, both absorption
lemmas, positivity of the five forms, the inversion) are unconditional
Lean theorems, while the literature inputs enter as \emph{named
hypotheses} of the main theorems --- \texttt{hCW} (the Corteel--Welsh
bounded forms \cite[Theorem~1.2]{CW19}, with denominators cleared),
\texttt{hQ} (the $Q$-transform relation of Theorem~\ref{thm:Q}, whose
Euler-series proof is not formalized), and, for the $d=8$ chain,
\texttt{hWarnaar} (Step~1), \texttt{hSplit} (the series form of Step~2,
whose polynomial kernel \emph{is} formalized as \texttt{qpoch\_split}),
\texttt{hBridge} (Step~3, \cite[Lemma~9.2]{KR22}) and \texttt{hUncu}
(Step~4, \cite[Theorem~4.3]{Uncu23}). The $d=10$ chain of
Theorem~\ref{thm:d10} has \emph{not} been formalized; its hypotheses
would be the exact analogues of the $d=8$ ones (with no bridge). A
machine-checked non-vacuity
witness shows that the hypotheses of the $d=4$ theorems are satisfiable.
Every theorem in the library depends only on the standard axioms
\texttt{propext}, \texttt{Classical.choice}, \texttt{Quot.sound}.

Separately from the formal proofs, Conjecture~\ref{con:pos} was verified
in exact $\ZZ[q]$ arithmetic (no series truncation), for \emph{every}
profile at each level, at
\[
d\in\{2,4,5,7,8,13,16,17,19,20,22,23,25,31\},
\]
in the ranges $n\le16$ for $d=13$; $n\le12$ for $d=16,17,19,20$;
$n\le10$ for $d=22,23$; $n\le9$ for $d=25$; $n\le7$ for $d=31$; and in
extended ranges at small level ($n\le25$ at $d=4$, $n\le22$ at $d=5$,
$n\le18$ at $d=7$, $n\le16$ at $d=8$, $n\le12$ at $d=13$, wall orbits
included). The minimal coefficients show no shrinking-margin trend as
$d$ grows.

\subsection*{Acknowledgements}
Everything here is built on the work of Corteel--Welsh,
Corteel--Dousse--Uncu, Kanade--Russell, Uncu, Welsh and Warnaar; the
contribution of this note is only to observe that their results, chained
together, prove a little more than has so far been recorded.

\appendix

\section{Methodology}\label{app:method}

The results of this note were produced by an orchestrated system of AI
agents (Claude, Anthropic), and the author believes it is right to say so
plainly. This appendix records how the system worked and what was done to
keep it honest. The author has read the mathematics of this note in full
and vouches for it; the division of labour was that the AI system
generated and checked the proofs, while the human role was that of
guarantor.

\subsection*{Corpus and seeding}
The system worked over a retrieval corpus of $82$ papers on
Rogers--Ramanujan-type identities, cylindric partitions and
neighbouring topics, split into roughly $7{,}000$ chunks and embedded with a fine-tuned SPECTER2
model. Eight initial research directions were selected by $k$-medoids
clustering of the embedded corpus, so that the starting points were
spread across genuinely different regions of the literature (bounded
forms, contiguous relations, vertex operators, bijective methods, and so
on) rather than chosen by hand.

\subsection*{The loop}
Each round (``layer'') ran eight agents in parallel, one per research
direction, each attempting proof strategies and recording partial
progress, failures and computations in a working file. A ninth agent then
synthesized the eight working files into a single document, and the next
layer was re-seeded from that synthesis. Later rounds allowed the agents
to generate their own retrieval queries against the corpus as their
attempts developed. The chain of results in this note --- the
$Q$-transform equivalence, the four-step $d=8$ assembly, and the $d=4$
absorption lemmas --- emerged over several such layers.

\subsection*{Discipline}
Since language models can produce confident errors, the system was built
around redundancy and adversarial checking rather than trust:
\begin{itemize}
\item One seed was a dedicated \emph{adversary}, tasked not with proving
anything but with independently recomputing others' claims and hunting
for counterexamples.
\item Any claim of a completed proof triggered independent verifier
agents, which re-derived the argument from the primary sources without
access to the claimant's reasoning.
\item The load-bearing elementary identities were formalized in
Lean~4/Mathlib (Section~\ref{sec:lean}; the repository
\url{https://github.com/RaggedR/warnaar-glue} is public), with the
literature inputs isolated as named hypotheses so that the boundary
between what is machine-checked and what is cited is explicit.
\item All positivity claims were tested by exact $\ZZ[q]$ computation
against an engine that had itself been validated against the raw
combinatorial definition of cylindric partitions, in the ranges reported
in Section~\ref{sec:lean}.
\end{itemize}
This redundancy earned its keep: two convention errors (a mislabelling of
profile orbits and an orientation error in a contiguous relation) were
caught precisely because multiple agents re-derived the same objects
independently and their answers disagreed; both were traced and fixed
against the primary sources. The standing rule that any bridging
contiguous relation be re-derived directly from \cite[Lemma~9.2]{KR22},
rather than taken from a restatement, is a lesson from one of these
episodes.

\subsection*{What this does and does not claim}
No claim is made that the system displayed insight beyond assembling,
checking and slightly extending results already in the literature; as
stated in the acknowledgements, the hard analysis is due to the cited
authors. The point of recording the methodology is transparency: the
reader should know that the proofs were found and verified this way, and
should weigh the machine-checked and exactly-computed components
(Section~\ref{sec:lean}) accordingly.

\begin{thebibliography}{99}

\bibitem{ASW99}
G.~E. Andrews, A.~Schilling and S.~O. Warnaar,
\textit{An $A_2$ Bailey lemma and Rogers--Ramanujan-type identities},
J. Amer. Math. Soc. \textbf{12} (1999), 677--702.

\bibitem{CDU22}
S.~Corteel, J.~Dousse and A.~K. Uncu,
\textit{Cylindric partitions and some new $A_2$ Rogers--Ramanujan
identities},
Proc. Amer. Math. Soc. \textbf{150} (2022), 481--497.

\bibitem{CW19}
S.~Corteel and T.~A. Welsh,
\textit{The $A_2$ Rogers--Ramanujan identities revisited},
Ann. Comb. \textbf{23} (2019), 683--694.

\bibitem{KR22}
S.~Kanade and M.~C. Russell,
\textit{Completing the $A_2$ Andrews--Schilling--Warnaar identities},
Int. Math. Res. Not. IMRN \textbf{2023}, no.~20, 17100--17155;
arXiv:2203.05690.

\bibitem{Comp26}
R.~Langer,
\textit{Warnaar's cylindric conjecture for the balanced profile
$(4,3,3)$ at modulus $13$: a complete proof via Warnaar's finite form
($a=+1$) and Uncu's modulo-$13$ theorem},
companion note to the present one, July 2026.

\bibitem{Uncu23}
A.~K. Uncu,
\textit{Proofs of modulo 11 and 13 cylindric Kanade--Russell conjectures
for $A_2$ Rogers--Ramanujan type identities},
arXiv:2301.01359.

\bibitem{W23}
S.~O. Warnaar,
\textit{The $\mathrm{A}_2$ Andrews--Gordon identities and cylindric
partitions},
Trans. Amer. Math. Soc. Ser. B \textbf{10} (2023), 715--765;
arXiv:2111.07550.

\bibitem{Welsh21}
T.~A. Welsh, unpublished.

\end{thebibliography}

\end{document}
